\subsection{Ideation und Lösungsraum}

Basierend auf den Erkenntnissen aus der Define Phase und den How-might-we-Fragen wurden Lösungsideen generiert, um die identifizierten Barrieren zu adressieren. Der Massnahmenraum konzentriert sich auf Kommunikation, Services, Erlebnisse und persönliche Begleitung --- den vom ZVV vorgegebenen Handlungsraum.

Der Innovationsprozess folgte dabei dem Principles of Design Thinking: Vielzahl von Ideen generieren, keine vorschnelle Bewertung, auf den Erkenntnissen der Forschung aufbauen, iterativ mit den Nutzer:innen testen.

\subsection{Drei Gestaltungsprinzipien}

Die Entwicklung von Massnahmen folgt drei übergeordneten Gestaltungsprinzipien:

\subsubsection{1. Narrativ-Shift: Auto vs. ÖV}

Das zentrale psychologische Problem ist nicht die Funktion des öV, sondern das Narrativ, das damit verbunden ist. Das Auto steht für Freiheit, Spontanität, Kontrolle, Unabhängigkeit. Der öV steht in der Wahrnehmung für Abhängigkeit, Fremdbestimmung, Einschränkung.

Gestaltungsprinzip: Umpositionierung des öV von ``notwendiges Verkehrsmittel'' (funktional) zu ``Möglichkeit für Erlebnis, Entspannung und Spontanität'' (emotionaler Benefit).

Konkret: Nicht ``Fahren Sie mit dem öV zur Kostenersparnis,'' sondern ``Fahren Sie mit dem öV und geniessen Sie Ihre Zeit --- lesen, entspannen, unterhalten.''

\subsubsection{2. Partnerschaftsmodell: Vertrauensaktivierung}

Die Stakeholder-Analyse zeigte: Der ZVV allein erreicht diese Zielgruppe nicht. Organisationen mit bestehendem Vertrauen (Pro Senectute, Gemeinden, Seniorenverbände, Apotheken, Hausärzte) spielen eine entscheidende Rolle.

Gestaltungsprinzip: Integration von Vertrauensorganisationen als Co-Aktivatoren. Diese Organisationen haben regelmässigen, niederschwelligen Zugang zur Zielgruppe und hohes Vertrauen. Sie können als Multiplikatoren fungieren.

\subsubsection{3. Transitions-Design: Führerscheinabgabe als Fenster}

Der kritische Moment der Führerscheinabgabe ist das stärkste bekannte Aktivierungsfenster für den öV. Bei früher Vorbereitung --- mindestens 1 Jahr vor der Abgabe --- ist die spätere öV-Nutzungsrate signifikant höher.

Gestaltungsprinzip: Spezielle Fokussierung auf den kritischen Moment als Aktivierungskanal. Massnahmen sollten optimal etwa 1 Jahr vor der möglichen Führerscheinabgabe (Fahrtauglichkeitsuntersuchung ab 70 Jahren) ansetzen.

\subsection{Prototyping und Testing}

Entwickelte Massnahmen wurden in Form von Low-Fidelity Prototypen (Skizzen, Szenarien, Rollenspiele) und wo möglich High-Fidelity Prototypen (funktionsfähige Services, Kampagnen-Mockups) konkretisiert. Diese wurden iterativ mit Vertreter:innen der Zielgruppe und der Partnerorganisationen getestet.

Testobjekte includeten:
Die Testobjekte umfassten Service-Szenarien (z.B. Ausflugplanung für Gruppen), Kommunikationsmaterialien und Kampagnenkonzepte sowie digitale und analoge Zugangskanäle. Zusätzlich wurden unterschiedliche Kooperationsmodelle mit Partnerorganisationen geprüft, um Umsetzbarkeit und Wirkung im Alltag früh abzusichern.

Das Feedback wurde kontinuierlich gesammelt und führte zu Anpassungen der Massnahmen.

\subsection{Dimensionen der Lösungsideen}

Aus der Ideation entstanden Lösungsansätze in mehreren Dimensionen:

\textbf{Kommunikation:} Neue Kommunikationsbotschaften und -kanäle, die das Erlebnis-Narrativ des öV stärken. Kanäle: analoge Sensibilisierung (Gemeinden, Apotheken), Mund-zu-Mund (Botschafter wie Kurt), niederschwellige Information.

\textbf{Services:} Neue Service-Angebote, die konkrete Hürden senken (z.B. Gruppen-Ticketing, Führerscheinabgabe-Begleitung, ÖV-Schnuppertage mit Guided Tours).

\textbf{Erlebnisse:} Gestaltete Erlebnisse, die den emotionalen Benefit des öV erlebbar machen (z.B. Themencafés in öV-Fahrzeugen, ÖV-Ausflugspakete mit Freizeitanbietern).

\textbf{Persönliche Begleitung:} Persönliche Support-Angebote für Transitionen (z.B. Buddy-Systeme, ÖV-Navigations-Support für Erstnutzende).

\subsection{Messbarer Erfolg}

Für das Prototyping wurden bereits Erfolgskriterien definiert:
Im Mittelpunkt standen die Teilnehmendenzahl und die Regelmässigkeit der öV-Nutzung, die Zufriedenheit mit den getesteten Formaten, die Conversion von Schnupper- zu Dauerangeboten sowie Weitergabeeffekte über Mund-zu-Mund-Kommunikation durch Botschafter:innen. Ergänzend wurde die Qualität der Zusammenarbeit mit Partnerorganisationen als eigener Erfolgsfaktor erfasst.

\subsection{Konkrete Lösungsbeispiele}

Aus der Ideation entstanden mehrere konkrete Lösungskonzepte, die in der Pilot-Phase getestet werden:

\subsubsection{Beispiel 1: ``ÖV-Botschafter-Programm''}

Eine strukturierte Initiative, um Personen wie Kurt zu aktivieren und auszubilden:
Das Programm beginnt mit der Rekrutierung aktiver ÖV-Nutzer:innen über bestehende Abos, Vereinsstrukturen und Befragungen. Darauf folgt ein kompaktes Training zu Mobilitätsrecht, Systemsicherheit und dem Umgang mit unsicheren Erstnutzenden. In der Aktivierungsphase laden Botschafter:innen ihr Umfeld zu begleiteten Schnupperfahrten ein und erhalten dafür passendes Material. Kleine Anreize wie Rabatte oder Community-Events unterstützen die Motivation, während regelmässige Austauschformate den Wissenstransfer und die Qualität sichern.

Ziel: Erreiche 30--50 Botschafter pro Pilotgemeinde, jede/r aktiviert durchschnittlich 3--5 neue Nutzer:innen.

\subsubsection{Beispiel 2: ``Führerscheinabgabe-Begleitung''}

Ein systematisches Programm für den kritischen Moment der Führerscheinabgabe:
Das Programm setzt früh an, indem Informationsmaterial etwa ein Jahr vor einer möglichen Führerscheinabgabe über Hausärzt:innen und Behörden verteilt wird. Kernstück ist ein Coaching mit fünf bis sechs Kontakten: einem Einstiegskurs zum Systemverständnis, mehreren begleiteten Fahrten mit schrittweiser Selbstständigkeit, einer ersten eigenständigen Fahrt und einer strukturierten Nachbesprechung. Anschliessend werden die Teilnehmenden an bestehende öV-Netzwerke, Vereine und Gruppen angebunden, damit neue Routinen sozial stabilisiert werden.

Ziel: Von den Personen, die einen Führerschein abgeben, nutzen mindestens 50\% regelmässig den öV.

\subsubsection{Beispiel 3: ``Gruppen-Ticketing und ÖV-Ausflugspakete''}

Service-Verbesserung für die häufige Use Case von Freizeitfahrten in Gruppen:
Für Gruppenreisen wurde ein Serviceansatz entwickelt, der ein einfaches App- oder Web-Ticketing mit transparenten Tarifen und automatischer Best-Deal-Logik verbindet. Ergänzend sollen Kooperationen mit Freizeitanbietern vorgepackte Ausflugspakete ermöglichen, während Gruppenrabatte den Einstieg erleichtern und den wahrgenommenen Planungsaufwand reduzieren.

Ziel: Reduktion des Planungs-Pain-Points ``Ticketing für Gruppen''.

\subsubsection{Beispiel 4: ``Analoge Schnellstart-Kits für Offline-Senior:innen''}

Für die ca. 30\% Offline-Senior:innen, die nicht digital erreichbar sind:
Für nicht-digitale Zielgruppen setzt das Schnellstart-Kit auf physische Verteilungspunkte wie Apotheken, Gemeindezentren und Arztpraxen. Inhaltlich kombiniert es eine gedruckte Einstiegshilfe zum Ticketkauf mit Netzkarte und alltagsnahen Zielbeispielen samt Fahrplänen; ergänzt wird dies durch eine Hotline und Vor-Ort-Unterstützung über Partnerorganisationen.

Ziel: Niederschwelliger Zugang für nicht-digitale Zielgruppe.

\subsection{Testing-Ansätze und Iteration}

Die vier Lösungskonzepte wurden in niedrig- und hochfidelity Prototypen getestet:

\subsubsection{Low-Fidelity Testing: Szenarien und Rollenspiele}

Mit Vertreter:innen der Personas wurden Szenarien durchgespielt:

\textit{Szenario 1: Kurt möchte einen Ausflug für Gruppe organisieren}
Ausgangslage war ein geplanter Säntis-Ausflug mit sechs Vereinskollegen. Die Kette aus Inspiration (Vereinsmail), Koordination (Telefon), Planung (interaktive Routen- und Kostenübersicht) und Ticketing zeigte klar: Der grösste Pain Point liegt in der Ticketkomplexität für Gruppen. Eine One-Click-Lösung mit automatischer Tarifoptimierung wurde in allen Durchläufen als potenzieller Game-Changer bewertet.

\textit{Szenario 2: Maria (ländlich) plant Besuch im Kunstmuseum}
Bei Maria wurde sichtbar, dass nicht primär der Wille, sondern die wahrgenommene Unsicherheit bei Taktung und Verbindungslogik hemmt. Eine visuell geführte Informationsseite mit Kartenbild, Fahrplanpfad und konkreten Optionen reduzierte die Hürde deutlich stärker als textlastige Zoneninformationen.

\textit{Szenario 3: Heinz nach Führerscheinabgabe}
Im dritten Szenario stand Heinz nach Führerscheinabgabe vor seiner ersten selbst organisierten Zugfahrt zu den Enkeln. Das Zusammenspiel aus Buddy-Begleitung, kurzem telefonischem Vorab-Support, Erinnerung vor der Fahrt und Check-in danach zeigte, dass emotionale Sicherheit mindestens so wichtig ist wie reine Sachinformation.

\subsubsection{High-Fidelity Testing: Prototypen mit Partner:innen}

Mit ausgewählten Partnerorganisationen wurden Materialien und Services getestet:
Mit ausgewählten Partnerorganisationen testeten wir ein ÖV-Botschafter-Handbuch inklusive Trainingsunterlagen, Mockups für digitales Gruppen-Ticketing, gedruckte Schnellstart-Kits für Offline-Senior:innen sowie Workshop-Formate zur Partneraktivierung. Das Feedback fiel überwiegend positiv aus und führte vor allem zu Vereinfachungen in Sprache, Prozessschritten und Übergabepunkten.

Das Feedback war durchgehend positiv, mit wichtigen Anpassungen zur kulturellen Angepasstheit und zur Vereinfachung der Prozesse.
